\documentclass{notaaula}

        \titulonota{Física moderna}
        \creditos{Turma 3A · Física · Dezembro/2026 · Prof. Josué Cavalcante}

        \begin{document}
        \cabecalho

        \begin{sobre}
        Introdução à física quântica, à física nuclear e à relatividade, com limites e evidências dos modelos. Habilidades em foco: EM13CNT205 e EM13CNT303.
        \end{sobre}

        \section{Quantização da energia}

\begin{copiar}{Fótons}
A luz troca energia em pacotes chamados fótons. A energia de um fóton é
\[
  E=hf=\frac{hc}{\lambda},
\]
com $h=\dec{6,63}\times10^{-34}\un{J\,s}$. Frequência maior corresponde a fóton mais energético;
comprimento de onda maior corresponde a energia menor.
\simbolos{$E$ (J); $f$ (Hz); $\lambda$ (m); $c=3\times10^8\un{m/s}$.}
\end{copiar}

\begin{exemplo}
Para $f=6\times10^{14}\un{Hz}$,
\[
  E=hf=\dec{6,63}\times10^{-34}\cdot6\times10^{14}
  \approx\resultado{4\times10^{-19}\un{J}}.
\]
\end{exemplo}

\begin{copiar}{Efeito fotoelétrico}
Um metal emite elétrons quando recebe fótons com energia mínima igual à função trabalho $\phi$:
\[
  K_{\max}=hf-\phi.
\]
Abaixo da frequência de corte, aumentar apenas a intensidade não libera elétrons. Acima do corte,
a frequência controla a energia máxima; a intensidade controla principalmente a quantidade de fótons.
\end{copiar}

\section{Átomos e espectros}

\begin{copiar}{Níveis discretos}
No modelo de Bohr, elétrons ocupam níveis de energia permitidos. Uma transição emite ou absorve fóton:
\[
  \Delta E=hf.
\]
Cada elemento apresenta linhas espectrais características, usadas para identificar composição de
estrelas, gases e materiais. Modelos quânticos atuais refinam o modelo de Bohr, mas preservam a ideia
de níveis discretos.
\end{copiar}

\section{Núcleo e radioatividade}

\begin{copiar}{Decaimento e meia-vida}
Núcleos instáveis podem emitir radiações alfa, beta ou gama. O número de núcleos não decaídos após
um tempo $t$ é
\[
  N=N_0\left(\frac12\right)^{t/T_{1/2}}.
\]
A meia-vida $T_{1/2}$ é o tempo para a quantidade cair à metade. O processo é probabilístico para
cada núcleo, mas previsível estatisticamente para uma amostra grande.
\end{copiar}

\begin{exemplo}[Três meias-vidas]
Uma amostra começa com $80\un{mg}$ e tem meia-vida de 4 dias. Após 12 dias,
\[
  80\rightarrow40\rightarrow20\rightarrow
  \resultado{10\un{mg}}.
\]
\end{exemplo}

\begin{copiar}{Fissão e fusão}
Fissão divide um núcleo pesado e é usada em reatores atuais. Fusão une núcleos leves, alimenta as
estrelas e exige condições extremas. Em ambos os casos, uma pequena variação de massa pode aparecer
como energia:
\[
  E=\Delta m\,c^2.
\]
Segurança nuclear exige controle de radiação, calor e resíduos; risco e benefício dependem de dose,
tecnologia e contexto.
\end{copiar}

\section{Relatividade restrita}

\begin{copiar}{Espaço e tempo}
As leis da Física têm a mesma forma em referenciais inerciais, e a rapidez da luz no vácuo é a mesma
para todos eles. Em velocidades próximas a $c$, intervalos de tempo e comprimentos dependem do
referencial. O fator
\[
  \gamma=\frac{1}{\sqrt{1-v^2/c^2}}
\]
cresce com $v$. Esses efeitos são desprezíveis em velocidades cotidianas, mas mensuráveis em
partículas rápidas e relevantes para tecnologias de navegação por satélite.
\end{copiar}

\begin{dica}
Física clássica não fica ``errada'': ela é uma excelente aproximação para objetos macroscópicos em
velocidades muito menores que $c$. Modelos têm domínios de validade.
\end{dica}

\begin{exercicios}
\nivel{1}{Conceitos}
\begin{questoes}
\item Como a energia de um fóton varia com a frequência?
\item O que é a frequência de corte no efeito fotoelétrico?
\item O que produz uma linha espectral no modelo de níveis de energia?
\item Defina meia-vida.
\item Diferencie fissão de fusão.
\nivel{2}{Aplicação}
\item Se a frequência de um fóton dobra, o que acontece com sua energia?
\item Uma amostra de $160\un{g}$ passa por duas meias-vidas. Quanto resta?
\item Uma substância tem meia-vida de 5 anos. Que fração resta após 15 anos?
\item Calcule a energia associada a $1\times10^{-9}\un{kg}$ usando $c=3\times10^8\un{m/s}$.
\nivel{3}{Síntese}
\item Explique por que intensidade alta não compensa frequência abaixo do corte fotoelétrico.
\item Compare previsibilidade de um núcleo individual e de uma grande amostra radioativa.
\item Explique o sentido de ``domínio de validade'' ao comparar mecânica clássica e relatividade.
\end{questoes}
\gabarito{1) cresce proporcionalmente; 6) dobra; 7) $40\un{g}$; 8) $1/8$;
9) $9\times10^7\un{J}$.}
\end{exercicios}


\end{document}
